\section{The one place finite precision changes a conclusion}
\label{sec:arithmetic}

Rounding perturbs every quantity above, mostly harmlessly. It matters in one place,
and that place earns a section because the conclusion it threatens is one of the
paper's claims: Proposition~\ref{prop:farkas} needs the entering constraint to be
genuinely violated, and in floating point that hypothesis is not free.

Notation for the working precision first. The implementation obtains it by a
construction due to Powell --- double a small positive number until it perturbs $1$
when scaled by both $0.1$ and $0.2$ --- giving a small multiple $\epsilon_0$ of the
machine epsilon without assuming what the arithmetic is. A constraint in the
working set has slack exactly zero by construction, and computed as
$c_i\T x - b_i$ it does not: the error is inherited from $x$ rather than generated
by the dot product, and is of order $\epsilon_0\norm{c_i}\norm{x}$. Working-set
slacks are therefore forced to zero rather than computed, and any slack below
$\epsilon_0$ is snapped to zero. That snap admits no principled threshold ---
separating ``violated by rounding'' from ``violated but slightly'' is impossible,
a genuine violation being arbitrarily small --- but getting it wrong costs a wasted
outer iteration rather than a wrong answer, since the loop re-measures every slack
at the new iterate.

Proposition~\ref{prop:farkas} requires the entering constraint to be genuinely
violated. Drop that hypothesis and the conclusion fails outright: $z = 0$ with no
limiting multiplier says only that the iterate cannot move and no multiplier can
be reduced, which for a constraint that is in fact satisfied is not a
contradiction but a description of a point that is already fine.

This is reachable, and not merely in principle. The choice of orthogonal reduction
in Section~\ref{ssec:insert} is free by Proposition~\ref{prop:invariance} but not
free of rounding: a Householder reflection and a Givens chain leave the iterate at
slightly different points, so a constraint that one implementation leaves a few
multiples of $\epsilon_0$ \emph{inside} its boundary, the other can leave a few
multiples \emph{outside} --- on either side of the snap above. Reaching the stuck state with such a constraint selected, the solver
must neither enforce it (there is nothing to enforce) nor conclude infeasibility
(the conclusion does not follow).

What it does instead is set the constraint aside for the current iterate and let
the outer loop take the next candidate, the set-aside list being cleared whenever
the iterate moves --- since a violation measured against a stale iterate says
nothing. The test for ``indistinguishable from rounding'' is
\begin{equation}
  |s_{i_*}| \;\le\; \kappa\, \epsilon_0 \max\bigl(
    \norm{c_{i_*}} \norm{x}_\infty + |b_{i_*}|,\; 1 \bigr)
  \label{eq:noise}
\end{equation}
for a fixed $\kappa$. Unlike that snap, this threshold
\emph{is} principled, and for an instructive reason: it does not have to separate
rounding from a small violation, only rounding from \emph{provable
infeasibility}. Infeasibility is macroscopic --- its size is set by the geometry of
the constraints, not by the arithmetic --- so the two populations are separated by
a dozen orders of magnitude and any threshold between them works. The scale in
\eqref{eq:noise} is $\norm{c}\norm{x}$ rather than the size of the terms forming
the dot product, precisely because the error is inherited from $x$.
